% Part: lambda-calculus
% Chapter: church-rosser
% Section: parallel-beta-eta-reduction

\documentclass[../../../include/open-logic-section]{subfiles}

\begin{document}

\olfileid[fr]{lam}{cr}{pbe}
\olsection{Réduction~\texorpdfstring{$\beta\eta$}{βη} parallèle}

Nous étudions la réduction~$\beta\eta$ parallèle et sa
propriété de Church--Rosser. Les règles s'appliquent aux
termes considérés à $\alpha$-équivalence près. La substitution
sur ces classes est définie par le \olref[lam][syn][alp]{cor:sub}.

\begin{defn}[Réduction~$\beta\eta$ parallèle,
$\beredpar$]\ollabel{defn:beredpar}
La \emph{réduction~$\beta\eta$ parallèle}
($\beredpar$) sur les termes est définie
inductivement par les règles suivantes :
\begin{enumerate}
\item \ollabel{defn:beredpar1}
  $x \beredpar x$.
\item \ollabel{defn:beredpar2}
  Si $N \beredpar N'$, alors
  $\lambd[x][N] \beredpar \lambd[x][N']$.
\item \ollabel{defn:beredpar3}
  Si $P \beredpar P'$ et $Q \beredpar Q'$,
  alors $PQ \beredpar P'Q'$.
\item \ollabel{defn:beredpar4}
  Si $N \beredpar N'$ et $Q \beredpar Q'$,
  alors $(\lambd[x][N])Q \beredpar
  \Subst{N'}{Q'}{x}$.
\item \ollabel{defn:beredpar5}
  Si $N \beredpar N'$, alors
  $\lambd[x][Nx]\beredpar N'$,
  à condition que $x\notin\FV{N}$.
\end{enumerate}
\end{defn}

Si $M\beredpar M'$, alors $\FV{M'}\subseteq\FV{M}$.
On le vérifie par induction sur la dérivation. L'application
réunit les ensembles de variables libres et l'abstraction en
retire le lieur. Dans la règle de contraction~$\beta$, les
variables libres de la substitution sont contenues dans
$(\FV{N'}\setminus\{x\})\cup\FV{Q'}$. Dans la règle~$\eta$,
l'hypothèse $x\notin\FV{N}$ donne
$\FV{\lambd[x][Nx]}=\FV{N}$.

Les règles sont invariantes par permutation des noms de variables.
Cela se vérifie sur chaque règle. En particulier, appliquer la
même permutation aux trois arguments d'une substitution permute
le résultat, puisque la substitution remplace les occurrences
libres et évite la capture. On peut ainsi ouvrir une abstraction
avec un nom frais dans une dérivation. On échange le nom de son
lieur avec un nom absent de la dérivation et des autres termes
considérés. La permutation agit sur le corps dans la prémisse.
L'abstraction obtenue représente la même classe. Dans une
contraction~$\eta$, le terme dont le lieur est retiré conserve
sa classe puisque ce nom n'y est pas libre.

\begin{thm}\ollabel{thm:refl}
$M \beredpar M$.
\end{thm}

\begin{proof}
Exercice.
\end{proof}

\begin{prob}
Démontrer le \olref[lam][cr][pbe]{thm:refl}.
\end{prob}

\begin{defn}[Développement~$\beta\eta$ complet]
\ollabel{defn:becd}
Le \emph{développement~$\beta\eta$ complet}
$\becd{M}$ de~$M$ est défini par les clauses
suivantes. Lorsqu'un terme est à la fois
une abstraction et un rédex~$\eta$
admissible, la dernière clause prévaut
sur la clause générale d'abstraction :
\begin{align}
  \becd{x} &= x \ollabel{defn:becd1} \\
  \becd{(\lambd[x][N])}
    &= \lambd[x][\becd{N}] \ollabel{defn:becd2}\\
  \becd{(PQ)} &= \becd{P}\becd{Q}
    && \text{si $P$ n'est pas une abstraction $\lambd$}
    \ollabel{defn:becd3} \\
  \becd{((\lambd[x][N])Q)}
    &= \Subst{\becd{N}}{\becd{Q}}{x}
    \ollabel{defn:becd4} \\
  \becd{(\lambd[x][Nx])}
    &= \becd{N} \ollabel{defn:becd5}
    && \text{si $x\notin\FV{N}$}
\end{align}
\end{defn}

Tous les appels récursifs portent sur des sous-termes stricts.
Le choix de la clause dépend de la structure du terme et, pour
un rédex~$\eta$, du lien entre le lieur et l'argument du corps,
ainsi que de la condition sur les variables libres. Deux
représentants $\alpha$-équivalents ont le même arbre étiqueté
au sens du \olref[lam][syn][alp]{thm:sub}. Sous un même nom
frais, leurs corps correspondants sont $\alpha$-équivalents.
La définition est donc indépendante du représentant par
induction sur le terme, en utilisant la substitution sur
les classes dans la clause de contraction~$\beta$.
Les mêmes clauses donnent $\FV{\becd{M}}\subseteq\FV{M}$
par induction sur $M$.

\begin{lem}\ollabel{lem:comp}
Si $M \beredpar M'$ et $R \beredpar R'$,
alors $\Subst{M}{R}{y}\beredpar
\Subst{M'}{R'}{y}$.
\end{lem}

\begin{proof}
Raisonnons par induction sur la dérivation
de $M\beredpar M'$. Les quatre premiers
cas suivent la preuve du
\olref[pb]{lem:comp}. Les règles de variable, d'application,
d'abstraction et de contraction~$\beta$ y sont les mêmes.
Dans le dernier de ces cas, l'identité de composition
s'applique aux termes des prémisses avec le lieur choisi frais.

Si la dernière étape applique la règle
\olref{defn:beredpar5}, écrivons
$M=\lambd[x][Nx]$ et $M'=N'$, avec
$x\notin\FV{N}$ et $N\beredpar N'$.
Choisissons par conversion~$\alpha$ un
représentant dont $x$ est frais pour
$y$, $R$ et $R'$. L'hypothèse
d'induction donne
$\Subst{N}{R}{y}\beredpar
\Subst{N'}{R'}{y}$.
De plus, $x\notin\FV{\Subst{N}{R}{y}}$.
La règle \olref{defn:beredpar5}
donne donc
\[
  \lambd[x][\Subst{N}{R}{y}x]
  \beredpar \Subst{N'}{R'}{y},
\]
ce qui est l'objectif sous le
représentant frais choisi.
\end{proof}

\begin{lem}\ollabel{lem:cont}
Si $M \beredpar M'$, alors
$M'\beredpar\becd{M}$.
\end{lem}

\begin{proof}
Raisonnons par induction forte sur le nombre de nœuds du
terme de départ. Pour les applications, nous prouvons
d'abord l'assertion auxiliaire suivante :
\[
  P\beredpar P',\quad Q\beredpar Q'
  \quad\Longrightarrow\quad P'Q'\beredpar\becd{(PQ)}.
\]
Si $P$ n'est pas une abstraction, les hypothèses d'induction
sur $P$ et $Q$ donnent $P'\beredpar\becd{P}$ et
$Q'\beredpar\becd{Q}$. La règle d'application donne la
conclusion. Si $P=\lambd[x][N]$ et si la réduction de $P$
se termine par la règle d'abstraction, alors
$P'=\lambd[x][N']$ avec $N\beredpar N'$. Les hypothèses
d'induction sur $N$ et $Q$, puis la règle de
contraction~$\beta$, donnent
$P'Q'\beredpar\Subst{\becd{N}}{\becd{Q}}{x}=\becd{(PQ)}$.

Il reste, pour cette assertion, le cas où
$P=\lambd[x][Lx]$ avec $x\notin\FV{L}$ et où la
\olref{defn:beredpar5} donne $P'$ à partir de
$L\beredpar P'$. L'application $LQ$ est strictement
plus petite que $PQ$. L'assertion auxiliaire pour $LQ$
donne $P'Q'\beredpar\becd{(LQ)}$.
Il suffit de vérifier que
\[
  \becd{((\lambd[x][Lx])Q)}\aeq\becd{(LQ)}.
\]
Choisissons $x$ frais pour $Q$ également. Si $L$
n'est pas une abstraction, le premier membre s'obtient
en substituant $\becd{Q}$ à $x$ dans $\becd{L}x$,
ce qui donne $\becd{L}\becd{Q}$. Ici $x$ n'est pas
libre dans $\becd{L}$. Si $L=\lambd[z][U]$, prenons
$x$ et $z$ distincts et frais pour $Q$, et $x$ frais
pour $U$. Le premier membre s'obtient en substituant
$\becd{Q}$ à $x$ dans $\Subst{\becd{U}}{x}{z}$.
Puisque $x$ n'est pas libre dans $\becd{U}$, ces deux
substitutions reviennent à insérer $\becd{Q}$ aux
occurrences libres de $z$ de $\becd{U}$, à
$\alpha$-équivalence près. C'est le second membre,
par la clause de contraction~$\beta$.
L'assertion auxiliaire est ainsi établie pour $PQ$.

Considérons maintenant la dernière règle de
$M\beredpar M'$.
\begin{enumerate}
\item Si $M$ est une variable, alors $M'=M=\becd{M}$.
\item Si la dernière règle est celle d'application, on vient
  d'établir la conclusion. Si elle contracte un rédex~$\beta$,
  écrivons $M=(\lambd[x][N])Q$ et
  $M'=\Subst{N'}{Q'}{x}$. L'hypothèse d'induction donne
  $N'\beredpar\becd{N}$ et $Q'\beredpar\becd{Q}$.
  Le \olref{lem:comp} donne la conclusion.
\item Si la dernière règle contracte un rédex~$\eta$,
  alors $M=\lambd[x][Lx]$ et $L\beredpar M'$.
  L'hypothèse d'induction donne
  $M'\beredpar\becd{L}=\becd{M}$.
\item Si la dernière règle est celle d'abstraction, écrivons
  $M=\lambd[x][N]$ et $M'=\lambd[x][N']$, avec
  $N\beredpar N'$. Si $M$ n'est pas un rédex~$\eta$,
  l'hypothèse d'induction sur $N$ et la règle d'abstraction
  donnent la conclusion comme dans le \olref[pb]{lem:cont}.
  Supposons donc $N=Lx$ et $x\notin\FV{L}$.
  \begin{enumerate}
  \item Si $N\beredpar N'$ se termine par la règle
    d'application, $N'=L'x$ avec $L\beredpar L'$.
    L'hypothèse d'induction donne $L'\beredpar\becd{L}$.
    La préservation des variables libres donne
    $x\notin\FV{L'}$. La contraction~$\eta$ donne donc
    $\lambd[x][L'x]\beredpar\becd{L}=\becd{M}$.
  \item Si $N\beredpar N'$ se termine par la
    contraction~$\beta$, alors $L=\lambd[z][U]$ et
    $N'=\Subst{U'}{x}{z}$, avec $U\beredpar U'$.
    Choisissons $z\neq x$. Puisque $x\notin\FV{L}$,
    $x\notin\FV{U}$ et donc $x\notin\FV{U'}$.
    Ainsi, $\lambd[x][\Subst{U'}{x}{z}]$ est
    $\alpha$-équivalent à $\lambd[z][U']$.
    Cette dernière abstraction est un réduit parallèle
    de $L$. L'hypothèse d'induction appliquée à $L$
    donne le pas vers $\becd{L}=\becd{M}$.
  \end{enumerate}
\end{enumerate}
\end{proof}

\begin{thm}\ollabel{thm:cr}
$\beredpar$ vérifie la propriété de
Church--Rosser.
\end{thm}

\begin{proof}
Le \olref{lem:cont} donne à deux réduits parallèles de $M$
un même successeur, $\becd{M}$, en un pas parallèle.
C'est la propriété de Church--Rosser.
\end{proof}

\end{document}
